% Physical p.12 / Roman X. Carrier word breaks joined; one missing source space corrected.
\IllusieSetPhysicalPageSize{447bp}{698.6bp}
\begin{illusiefrontpage}{illusie:I:intro:p012}{}
\begin{center}\fontsize{10}{12}\selectfont X\end{center}
\vspace{12bp}
\phantomsection\hypertarget{illusie:I:intro:p012:p1}{}%
Au chapitre II, en utilisant la théorie des résolutions standard et la variante du
théorème de Whitehead indiquée plus haut, nous définissons le complexe cotangent
$L_{X/Y}$ d'un morphisme de topos annelés $X\longrightarrow Y$, et étudions ses
propriétés générales de fonctorialité et compatibilité au changement de base. Le
résultat principal qui généralise ([25] 5.1), est qu'un composé
$X\xrightarrow{f}Y\longrightarrow Z$ de morphismes de topos annelés donne naissance
à un triangle distingué (canonique et fonctoriel) dans la catégorie dérivée $D(X)$,
dit ``triangle de transitivité'' :

\phantomsection\hypertarget{illusie:I:intro:p012:d1}{}%
\IllusieTransitivityTriangle

\vspace{1bp}
\phantomsection\hypertarget{illusie:I:intro:p012:p2}{}%
Dans le chapitre III, nous établissons, pour un morphisme de topos annelés
$X\longrightarrow Y$ et un $\mathcal O_X$-Module $I$, un isomorphisme fonctoriel
canonique entre le groupe des classes à isomorphisme près de $Y$-extensions de $X$
par $I$ et le hyperext $\operatorname{Ext}^{1}(L_{X/Y},I)$. Ce résultat nous permet
de comparer $L_{X/Y}$ au complexe construit par Grothendieck dans [17] : ce dernier
est isomorphe dans $D(X)$ au tronqué de $L_{X/Y}$ obtenu en tuant la cohomologie en
degré $\leq -2$. Combinant ces résultats avec ceux du chapitre II, nous obtenons des
applications de la théorie à des questions de déformations de topos annelés et de
morphismes de topos annelés. Nous exprimons l'obstruction à la déformation plate au
premier ordre d'un topos annelé comme cup-produit avec une certaine classe
fondamentale définie en termes du triangle de transitivité décrit plus haut, classe
qui apparaît comme une généralisation naturelle du classique invariant de
Kodaira--Spencer (cf. [16]). Ce résultat résout affirmativement une conjecture de
Grothendieck, qui est à l'origine du présent travail. Le reste du chapitre est
consacré au calcul du complexe cotangent de quelques types de morphismes de schémas
(morphismes lisses, intersections complètes).

\vspace{4bp}
\phantomsection\hypertarget{illusie:I:intro:p012:p3}{}%
Utilisant une variante de la théorie du complexe cotangent pour les faisceaux
d'algèbres graduées, nous traitons, au chapitre IV, des questions de déformations de
Modules. L'obstruction à la déformation plate au premier ordre d'un Module est
exprimée comme cup-produit avec une certaine classe fondamentale, qui généralise
\end{illusiefrontpage}
